% !TeX root = ../main_skripsi.tex
\chapter{HASIL PENELITIAN DAN PEMBAHASAN}

\section{Analisis Deskripsi}

[Jelaskan perkembangan dan karakteristik data berdasarkan keluaran analisis
terverifikasi. Jangan menyalin data mentah langsung ke naskah.]

\section{Statistik Deskriptif}

[Sajikan jumlah observasi, nilai rata-rata atau median, sebaran, minimum,
maksimum, dan ukuran lain yang relevan. Semua angka wajib cocok dengan keluaran
analisis terverifikasi.]

\section{Hasil Penelitian}

[Sajikan hasil objektif mengikuti urutan pertanyaan penelitian atau hipotesis.
Pisahkan hasil dari interpretasi. Cantumkan ukuran efek, ketidakpastian, dan
diagnostik model bila relevan.]

\section{Pembahasan}

[Tafsirkan temuan memakai teori pada Bab II, bandingkan dengan studi terdahulu,
jelaskan mekanisme yang mungkin, dan nyatakan kontribusi tanpa melampaui bukti.]

\section{Keterbatasan Penelitian}

[Uraikan keterbatasan desain, sampel, pengukuran, data, analisis, serta dampaknya
terhadap validitas dan generalisasi.]
